\thispagestyle{plain}

\chapter{Feature Engineering}
\label{cp:fe}
\vspace{1\baselineskip}

Based on our Exploratory Data Analysis, we will now construct two distinct feature engineering pipelines. This approach is necessary because linear
models (e.g., Logistic Regression) and tree-based models (e.g., Random Forest) have different preprocessing requirements.

\section{Pipeline for Linear Models}
\begin{itemize}
    \item NaN Imputation: Replace missing values in workclass and occupation with a new 'Unknown' category, as decided in our EDA.
    \item Binning: Discretize the age and education-num features into categorical bins to capture their non-linear relationships with income.
    \item Encoding: Convert all categorical features (including the newly binned ones) into numerical format using One-Hot Encoding.
    \item Scaling: Standardize all numerical features to have a mean of 0 and a standard deviation of 1. This is crucial to prevent features with larger scales from dominating the model.
\end{itemize}

\section{Pipeline for Tree-Based Models}
\begin{itemize}
    \item NaN Imputation: Same as the linear pipeline (replace with 'Unknown').
    \item Encoding: Convert all categorical features using One-Hot Encoding.
    \item No Binning/Scaling: We will not apply binning or scaling. Tree models are not sensitive to the scale of features and are capable of capturing non-linear relationships on their own by finding optimal split points.
\end{itemize}

\section{Target Variable Encoding}

Finally, we encode the target variable y from categorical ('$\le$50K', '>50K') to numerical (0, 1), which is required by the models.


